% ---------------- ABSTRAK (maks. 300 kata, satu spasi, tanpa indentasi) --
% TODO: perbarui hasil setelah Bab IV selesai.
\frontchapter{ABSTRAK}
\begingroup\setstretch{1}\setlength{\parindent}{0pt}
Layanan administrasi akademik di Fakultas Sains dan Matematika (FSM) Universitas Diponegoro dijalankan berdasarkan Standar Operasional Prosedur (SOP), sehingga staf tata usaha menerima banyak pertanyaan berulang mengenai urutan langkah, pejabat yang bertanggung jawab, dan dokumen persyaratan. Asisten virtual berbasis \textit{Large Language Model} (LLM) dapat menjawab pertanyaan tersebut, tetapi LLM berbasis \textit{cloud} menimbulkan biaya per kueri dan risiko kebocoran data mahasiswa, sedangkan \textit{Retrieval-Augmented Generation} (RAG) konvensional cenderung mencampurkan langkah dari SOP yang berbeda dan membalik urutan langkah. Penelitian ini mengembangkan asisten virtual yang terintegrasi dengan sistem persuratan digital E-Office dan menggunakan mesin \textit{Process-Aware Workflow} GraphRAG. Setiap SOP dimodelkan sebagai graf alur kerja berarah yang menyimpan aktor, aksi, dokumen masukan, keluaran, dan durasi setiap langkah. Pemilih SOP berbasis aturan menentukan prosedur yang relevan, langkah-langkahnya diserialisasi secara berurutan, lalu dilengkapi potongan dokumen pendukung sebelum dijawab oleh LLM lokal \texttt{gemma4:e2b} melalui Ollama. Alur percakapan diorkestrasi menggunakan LangGraph dan disajikan melalui antarmuka E-Office berbasis Next.js dengan layanan \textit{backend} Bun dan Elysia. Evaluasi menggunakan 40 pertanyaan terverifikasi dengan metrik \textit{key-fact recall} (KFR), cakupan langkah, \textit{Topological Inversion Rate} (TIR), \textit{Actor Mismatch Rate} (AMR), akurasi penolakan, dan efisiensi token, dilengkapi pengujian \textit{black-box} terhadap fitur sistem. Pengujian pendahuluan menunjukkan bahwa mesin tersebut mencapai KFR 0,955 dibandingkan 0,712 pada Naive RAG, cakupan langkah 1,000, TIR 0,000, dan menggunakan 44\% lebih sedikit token prompt dibandingkan pemberian seluruh dokumen. Hasil ini menunjukkan bahwa pemilihan prosedur dan representasi alur kerja yang eksplisit memungkinkan LLM berukuran kecil yang berjalan secara lokal memberikan jawaban prosedural yang lengkap dan berurutan.

\vspace{12pt}
\noindent Kata kunci : \textit{Retrieval-Augmented Generation}, GraphRAG, graf alur kerja, SOP, E-Office, asisten virtual, LLM lokal
\par\endgroup
\clearpage
